\documentclass[12pt, a4paper]{article}
\usepackage{amsmath, amssymb, amsthm}
\usepackage{booktabs}
\usepackage{geometry}
\usepackage{hyperref}
\usepackage{array}

\geometry{margin=2.5cm}

\newtheorem{theorem}{Theorem}[section]
\newtheorem{corollary}[theorem]{Corollary}
\theoremstyle{definition}
\newtheorem{definition}[theorem]{Definition}
\newtheorem{example}[theorem]{Example}
\theoremstyle{remark}
\newtheorem*{remark}{Remark}

\title{\textbf{Series and Integral Transforms}\\[0.5em]
\large Theorems \& Formulas Reference Sheet}
\author{Course 21183 --- HIT, Prof.\ Michael Kroyter}
\date{}

\begin{document}
\maketitle
\tableofcontents
\newpage

%----------------------------------------------------------------------
\section{Numerical Series}

\subsection{Basic Definitions and Facts}
\begin{itemize}
  \item The series $\displaystyle\sum_{n=1}^\infty a_n$ \textbf{converges} if the partial sums
        $S_N = \sum_{n=1}^N a_n$ have a finite limit.
  \item \textbf{Necessary condition}: if $\sum a_n$ converges, then $a_n \to 0$.
        Contrapositive: if $a_n \not\to 0$, the series diverges.
  \item \textbf{Absolutely convergent}: $\sum |a_n| < \infty$ (implies convergent).
  \item \textbf{Conditionally convergent}: $\sum a_n$ converges but $\sum |a_n|$ diverges.
  \item \textbf{$p$-series}: $\displaystyle\sum_{n=1}^\infty \frac{1}{n^p}$ converges $\Leftrightarrow$ $p > 1$.
\end{itemize}

\subsection{Convergence Tests}

\begin{center}
\begin{tabular}{>{\bfseries}p{2.8cm}p{4.5cm}p{5.5cm}}
\toprule
Test & Preconditions & Conclusion \\
\midrule
Comparison
  & $a_n, b_n \ge 0$ \newline (positive series)
  & If $a_n \le b_n$ and $\sum b_n<\infty$, then $\sum a_n <\infty$.
    If $a_n \ge b_n \ge 0$ and $\sum b_n=\infty$, then $\sum a_n=\infty$. \\[6pt]
Limit Comparison
  & $a_n, b_n > 0$ \newline (positive series)
  & If $\lim\dfrac{a_n}{b_n} = L \in (0,\infty)$, then $\sum a_n$ and $\sum b_n$ behave the same. \\[6pt]
Ratio (D'Alembert)
  & $a_n \ne 0$; tests \emph{absolute} convergence
  & $L = \lim\!\left|\dfrac{a_{n+1}}{a_n}\right|$: converges if $L<1$, diverges if $L>1$, inconclusive if $L=1$. \\[6pt]
Root (Cauchy)
  & Any series; tests \emph{absolute} convergence
  & $L = \limsup|a_n|^{1/n}$: converges if $L<1$, diverges if $L>1$. \\[6pt]
Integral
  & $f$ continuous, positive, \emph{decreasing} on $[1,\infty)$;\newline $a_n = f(n)$\newline (conditions need hold only eventually)
  & $\sum_{n=1}^\infty f(n)$ and $\int_1^\infty\! f(x)\,dx$ converge or diverge together. \\[6pt]
Alternating (Leibniz)
  & $b_n \ge 0$;\newline $b_n \searrow 0$ monotonically
  & $\sum(-1)^n b_n$ converges. \\[6pt]
Dirichlet
  & Partial sums of $\sum a_n$ bounded;\newline $b_n \searrow 0$ monotonically
  & $\sum a_n b_n$ converges. \\[6pt]
Abel
  & $\sum a_n$ converges;\newline $b_n$ monotone and bounded
  & $\sum a_n b_n$ converges. \\
\bottomrule
\end{tabular}
\end{center}

\begin{remark}[Parity trick: $\cos(n^k\pi) = (-1)^n$ for odd $k$]
For any integer $n$ and \emph{odd} positive integer $k$: $n^k \equiv n \pmod{2}$, so
$\cos(n^k\pi) = (-1)^{n^k} = (-1)^n$.
Similarly $\sin(n^k\pi) = 0$ for any positive integer $k$.
Use this to simplify series whose terms contain $\cos(n^3\pi)$, $\cos(n^5\pi)$, etc.
\end{remark}

\begin{theorem}[Alternating Series Error Bound]
If $\sum(-1)^n b_n$ satisfies the Leibniz conditions ($b_n \ge 0$, $b_n \searrow 0$), then
the error after $N$ terms satisfies $|S - S_N| \le b_{N+1}$.
To achieve accuracy $<\varepsilon$: use the smallest $N$ with $b_{N+1} < \varepsilon$.
\end{theorem}

\begin{theorem}[Telescoping Series]
If $a_n \to L$ (finite), then $\displaystyle\sum_{n=1}^\infty (a_n - a_{n+1}) = a_1 - L$,
since $S_N = a_1 - a_{N+1} \to a_1 - L$.
Recognition: a series written as $f(n) - f(n+1)$ (or $f(n+1) - f(n)$) is telescoping.
\end{theorem}

%----------------------------------------------------------------------
\section{Power Series}

\subsection{Radius of Convergence}
The series $\displaystyle\sum_{n=0}^\infty c_n (x-a)^n$ converges absolutely for $|x-a| < R$, where
\[
  R = \frac{1}{\displaystyle\limsup_{n\to\infty} |c_n|^{1/n}}
  \qquad \text{(Cauchy--Hadamard)}
  \qquad \text{or} \qquad
  R = \lim_{n\to\infty} \left|\frac{c_n}{c_{n+1}}\right| \text{ when the limit exists.}
\]
The series diverges for $|x-a|>R$; endpoints $|x-a|=R$ must be checked separately.
Integration and differentiation do not change $R$.

\subsection{Taylor and Maclaurin Series}
\[
  f(x) = \sum_{n=0}^\infty \frac{f^{(n)}(a)}{n!}(x-a)^n
\]

\textbf{Reading off derivatives:} if $f(x) = \displaystyle\sum_{n=0}^\infty c_n x^n$, then
\[
  \boxed{f^{(k)}(0) = k!\, c_k}
\]
To find $f^{(k)}(0)$: read off the coefficient of $x^k$ in the series and multiply by $k!$.
No need to differentiate the series $k$ times explicitly.

\subsection{Product and Quotient of Power Series}

\textbf{Cauchy product:} if $f(x) = \displaystyle\sum_{n=0}^\infty b_n x^n$ and
$g(x) = \displaystyle\sum_{n=0}^\infty c_n x^n$, then
\[
  f(x)\,g(x) = \sum_{n=0}^\infty a_n x^n, \qquad a_n = \sum_{k=0}^n b_k\,c_{n-k}.
\]
The radius of convergence of the product is at least $\min(R_f, R_g)$.

\textbf{Example:} $\dfrac{-\ln(1-x)}{1-x}
  = \!\left(\sum_{n=1}^\infty \tfrac{x^n}{n}\right)\!\!\left(\sum_{n=0}^\infty x^n\right)$;
the coefficient of $x^n$ is $a_n = H_n = \sum_{k=1}^n \tfrac{1}{k}$ (harmonic number).

\textbf{Quotient of power series:} to expand $h(x) = P(x)/Q(x)$ near $x=0$ (with $Q(0)\ne 0$):
\begin{enumerate}
  \item Expand $P$ and $Q$ as power series to the required order.
  \item Write $Q(x) = Q(0)\bigl(1+r(x)\bigr)$, $r(0)=0$; use $1/(1+r) = 1-r+r^2-\cdots$
  \item Multiply $P(x)\cdot\tfrac{1}{Q(x)}$ term by term, then read off $h^{(k)}(0) = k!\cdot[\text{coeff of }x^k]$.
\end{enumerate}

\subsection{Common Maclaurin Series}
\begin{align*}
  e^x &= \sum_{n=0}^\infty \frac{x^n}{n!}, \qquad R = \infty \\[4pt]
  \sin x &= \sum_{n=0}^\infty \frac{(-1)^n\, x^{2n+1}}{(2n+1)!}, \qquad R = \infty \\[4pt]
  \cos x &= \sum_{n=0}^\infty \frac{(-1)^n\, x^{2n}}{(2n)!}, \qquad R = \infty \\[4pt]
  \sinh x &= \sum_{n=0}^\infty \frac{x^{2n+1}}{(2n+1)!}, \qquad R = \infty
  \quad \text{(same as }\sin x\text{ without }(-1)^n\text{)} \\[4pt]
  \cosh x &= \sum_{n=0}^\infty \frac{x^{2n}}{(2n)!}, \qquad R = \infty
  \quad \text{(same as }\cos x\text{ without }(-1)^n\text{)} \\[4pt]
  \frac{1}{1-x} &= \sum_{n=0}^\infty x^n, \qquad |x| < 1 \\[4pt]
  \ln(1+x) &= \sum_{n=1}^\infty \frac{(-1)^{n+1}\, x^n}{n}, \qquad |x| \le 1,\ x \ne -1 \\[4pt]
  \arctan x &= \sum_{n=0}^\infty \frac{(-1)^n\, x^{2n+1}}{2n+1}, \qquad |x| \le 1 \\[4pt]
  (1+x)^\alpha &= \sum_{n=0}^\infty \binom{\alpha}{n} x^n, \qquad |x| < 1
\end{align*}

\textbf{Binomial coefficient reminder:}
\[
  \binom{n}{k} = \frac{n!}{k!\,(n-k)!}
  \qquad \text{(integer }n \ge k \ge 0\text{, ``}n\text{ choose }k\text{'')}
\]
\[
  \binom{\alpha}{k} = \frac{\alpha\,(\alpha-1)\cdots(\alpha-k+1)}{k!}
  \qquad \text{(generalized, any }\alpha \in \mathbb{R}\text{)}
\]
Pascal's rule: $\dbinom{n}{k} = \dbinom{n-1}{k-1} + \dbinom{n-1}{k}$.\quad
Special values: $\dbinom{n}{0} = \dbinom{n}{n} = 1$, $\dbinom{n}{1} = n$.

\subsection{Taylor's Theorem with Remainder}

\begin{theorem}[Lagrange Remainder]
If $f$ has $n+1$ continuous derivatives on an interval containing $a$, then
\[
  f(x) = \sum_{k=0}^n \frac{f^{(k)}(a)}{k!}(x-a)^k + R_n(x),
  \qquad
  R_n(x) = \frac{f^{(n+1)}(c)}{(n+1)!}(x-a)^{n+1}
\]
for some $c$ strictly between $a$ and $x$.  Error bound:
\[
  |R_n(x)| \leq \frac{M_{n+1}}{(n+1)!}\,|x-a|^{n+1},
  \qquad M_{n+1} = \sup_{t \text{ between } a,x}\bigl|f^{(n+1)}(t)\bigr|.
\]
\end{theorem}

\subsection{Applications: Limits and Integrals via Taylor}

\textbf{Computing limits at 0:} expand numerator and denominator as power series, cancel common $x^k$ factors, read off the limit from the leading terms.

\textbf{Computing integrals:} substitute a known series for the integrand and integrate term by term (valid inside the radius of convergence). Use the Leibniz error bound (Theorem 1.2) to determine how many terms are needed.

%----------------------------------------------------------------------
\section{Uniform Convergence}

\subsection{Sequences of Functions}

\begin{definition}
Let $f_n, f : D \to \mathbb{R}$.
\begin{itemize}
  \item \textbf{Pointwise}: $f_n \to f$ if for each $x \in D$, $f_n(x) \to f(x)$.
  \item \textbf{Uniform} ($f_n \rightrightarrows f$): $\displaystyle\sup_{x \in D}|f_n(x) - f(x)| \to 0$.
\end{itemize}
Uniform $\Rightarrow$ pointwise, but pointwise $\not\Rightarrow$ uniform in general.
\end{definition}

\textbf{Standard method:} compute $M_n = \sup_{x\in D}|f_n(x)-f(x)|$ (find the max by differentiation).
If $M_n \to 0$: uniform. If $M_n \not\to 0$: not uniform.

\textbf{Negation:} $f_n \not\rightrightarrows f$ iff $\exists\,\varepsilon > 0$ and a sequence $x_n \in D$ with $|f_n(x_n) - f(x_n)| \geq \varepsilon$.
Typical choice: $x_n = \arg\max_x |f_n(x) - f(x)|$.

\begin{example}[$f_n(x) = x^n$ on $[0,1)$]
Pointwise limit: $f(x) = 0$. But $\sup_{x\in[0,1)}|x^n| = 1 \not\to 0$. Not uniform.
On $[0,r]$ for any $r < 1$: $\sup = r^n \to 0$. Uniform.
\end{example}

\subsection{Properties Preserved by Uniform Convergence}

\begin{theorem}[Continuity]
If each $f_n$ is continuous on $D$ and $f_n \rightrightarrows f$, then $f$ is continuous on $D$.
\end{theorem}

\begin{theorem}[Term-by-term Integration]
If $f_n \rightrightarrows f$ on $[a,b]$, then
$\displaystyle\lim_{n\to\infty}\int_a^b f_n(x)\,dx = \int_a^b f(x)\,dx$.
\end{theorem}

\begin{theorem}[Term-by-term Differentiation]
If $f_n \to f$ pointwise on $[a,b]$, each $f_n$ is $C^1$, and $f_n' \rightrightarrows g$ uniformly,
then $f$ is differentiable and $f' = g$.
\end{theorem}

\begin{theorem}[Sequential Substitution]
If $f_n \rightrightarrows f$ uniformly on $[a,b]$ and $x_n \to x \in [a,b]$, then $f_n(x_n) \to f(x)$.
\end{theorem}

\begin{remark}[Swapping $\sum$ and improper integrals]
For $\int_0^\infty \sum u_n(x)\,dx$: if each $u_n \ge 0$ and $\sum \int_0^\infty u_n < \infty$,
the swap is valid. Otherwise: (1) justify the swap on each compact $[a,b]$ via uniform convergence;
(2) show the endpoint contributions (near $0$ and near $\infty$) vanish as $a\to 0^+$, $b\to\infty$.
\end{remark}

\subsection{Series of Functions}

A series $\displaystyle\sum_{n=1}^\infty u_n(x)$ \textbf{converges uniformly} on $D$ if its partial sums converge uniformly on $D$.

\textbf{Domain of convergence:} treat $x$ as a parameter; apply standard numerical-series tests to $\sum u_n(x)$ for each fixed $x$.

\begin{theorem}[Weierstrass $M$-Test]
If $|u_n(x)| \leq M_n$ for all $x \in D$ and $\displaystyle\sum M_n < \infty$, then
$\sum u_n(x)$ converges uniformly and absolutely on $D$.
\end{theorem}

\begin{remark}
The $M$-test is sufficient but \emph{not necessary}: a series can converge uniformly even when $\sum M_n = \infty$.
When the $M$-test fails, try the uniform Dirichlet or Abel tests.

\textbf{Example} ($M$-test fails, Dirichlet works): $\displaystyle\sum_{n=1}^\infty \frac{(-1)^n}{n}$
write as $\sum a_n b_n$ with $a_n = (-1)^n$, $b_n = 1/n$.
\begin{itemize}
  \item Partial sums of $\sum(-1)^n$: $|A_N| \le 1$ for all $N$ \quad (bounded).
  \item $b_n = 1/n \searrow 0$ monotonically (no $x$-dependence, so trivially uniform).
\end{itemize}
By the uniform Dirichlet test, $\sum(-1)^n/n$ converges uniformly on $\mathbb{R}$,
yet $\sum 1/n = \infty$ so the $M$-test cannot be applied.
\end{remark}

\begin{theorem}[Dirichlet --- Uniform Version]
$\sum a_n(x)\,b_n(x)$ converges uniformly on $D$ if the partial sums $A_N(x)$ are uniformly bounded
and $b_n(x) \searrow 0$ monotonically and uniformly on $D$.
\end{theorem}

\begin{theorem}[Abel --- Uniform Version]
$\sum a_n(x)\,b_n(x)$ converges uniformly on $D$ if $\sum a_n(x)$ converges uniformly
and $b_n(x)$ is monotone and uniformly bounded on $D$.
\end{theorem}

\textbf{Power series:} $\displaystyle\sum c_n x^n$ with radius $R$ converges uniformly on $[-r,r]$ for any $r < R$,
but generally not on $(-R,R)$.

\textbf{Properties of uniformly convergent series:} the sum is continuous if each $u_n$ is; term-by-term
integration is valid; term-by-term differentiation requires $\sum u_n'$ to converge uniformly.

%----------------------------------------------------------------------
\section{Inner Product Spaces \& $L^2$ Norm}

\subsection{Inner Product on $L^2[-L,L]$}

\[
  \langle f, g \rangle = \int_{-L}^{L} f(x)\,\overline{g(x)}\,dx,
  \qquad
  \|f\| = \sqrt{\langle f,f\rangle} = \sqrt{\int_{-L}^L |f(x)|^2\,dx}
\]

\subsection{Orthogonality of the Trigonometric Basis}

The functions $\{1,\,\cos(n\pi x/L),\,\sin(n\pi x/L)\}_{n\ge 1}$ are orthogonal on $[-L,L]$:
\[
  \int_{-L}^{L} \cos\frac{n\pi x}{L}\cos\frac{m\pi x}{L}\,dx
  = \begin{cases} 2L & n=m=0 \\ L & n=m\neq 0 \\ 0 & n\neq m \end{cases}
  \qquad
  \int_{-L}^{L} \sin\frac{n\pi x}{L}\sin\frac{m\pi x}{L}\,dx
  = \begin{cases} L & n=m\neq 0 \\ 0 & \text{otherwise} \end{cases}
\]
\[
  \int_{-L}^{L} \sin\frac{n\pi x}{L}\cos\frac{m\pi x}{L}\,dx = 0 \quad \text{for all } n,m
\]

\subsection{Bessel's Inequality and Parseval's Identity}

\begin{center}
\renewcommand{\arraystretch}{1.5}
\begin{tabular}{p{6.5cm}p{6.5cm}}
\toprule
\textbf{Bessel's Inequality} & \textbf{Parseval's Identity} \\
\midrule
$\displaystyle\frac{a_0^2}{2} + \sum_{n=1}^{N}(a_n^2+b_n^2) \le \frac{1}{L}\|f\|^2$
\newline for every finite $N$.
&
$\displaystyle\frac{a_0^2}{2} + \sum_{n=1}^{\infty}(a_n^2+b_n^2) = \frac{1}{L}\|f\|^2$
\newline (holds for all $f \in L^2[-L,L]$). \\[6pt]
Holds for any $f \in L^2[-L,L]$, even if the Fourier series does not converge to $f$ pointwise.
& Parseval = Bessel with equality.
Equivalently: $\|S_N - f\| \to 0$ (convergence in norm). \\
\bottomrule
\end{tabular}
\end{center}

\textbf{Convergence implications} (on a bounded interval):
uniform $\Rightarrow$ pointwise;\ uniform $\Rightarrow$ $L^2$ norm;\ pointwise $\not\Rightarrow$ $L^2$;\ $L^2 \not\Rightarrow$ pointwise.

%----------------------------------------------------------------------
\section{Fourier Series}

\subsection{Definition and Coefficients}
For $f$ periodic with period $2L$:
\[
  f(x) \sim \frac{a_0}{2} + \sum_{n=1}^\infty
  \left( a_n \cos\frac{n\pi x}{L} + b_n \sin\frac{n\pi x}{L} \right)
\]
\[
  a_n = \frac{1}{L}\int_{-L}^{L} f(x)\cos\frac{n\pi x}{L}\,dx,
  \qquad
  b_n = \frac{1}{L}\int_{-L}^{L} f(x)\sin\frac{n\pi x}{L}\,dx
\]
The integral may be taken over any full period, not only $[-L,L]$.

\subsection{Pointwise Convergence: Dirichlet's Theorem}
If $f$ is piecewise smooth and periodic, the Fourier series $S(x)$ converges to:
\begin{itemize}
  \item $f(x)$ at points of continuity,
  \item $\dfrac{f(x^+)+f(x^-)}{2}$ at jump discontinuities.
\end{itemize}

\textbf{Evaluating $S(x_0)$ outside $[-L,L]$:} The Fourier series represents the \emph{periodic extension}
of $f$ (period $2L$). For $x_0 \notin [-L,L]$: reduce $x_0$ modulo $2L$ to $x_0' \in [-L,L]$,
then apply Dirichlet's theorem at $x_0'$.

\textbf{Example:} $f$ on $[-\pi,\pi]$, find $S(3\pi)$.
Since $3\pi \equiv \pi \pmod{2\pi}$, and $\pi$ is a jump of the periodic extension:
$S(3\pi) = S(\pi) = \dfrac{f(\pi^-)+f(-\pi^+)}{2}$.

\subsection{Even, Odd, and Half-Period Extensions}

\textbf{Even $f$:} $b_n = 0$;\quad
$a_n = \dfrac{2}{L}\displaystyle\int_0^L f(x)\cos\dfrac{n\pi x}{L}\,dx$

\textbf{Odd $f$:} $a_n = 0$;\quad
$b_n = \dfrac{2}{L}\displaystyle\int_0^L f(x)\sin\dfrac{n\pi x}{L}\,dx$

\medskip
Given $f$ on $[0,L]$, extend to $[-L,L]$ (evenly or oddly) then periodically:

\textbf{Cosine series} (even extension):
\[
  f(x) \sim \frac{a_0}{2} + \sum_{n=1}^\infty a_n \cos\frac{n\pi x}{L},
  \qquad
  a_n = \frac{2}{L}\int_0^L f(x)\cos\frac{n\pi x}{L}\,dx
\]

\textbf{Sine series} (odd extension):
\[
  f(x) \sim \sum_{n=1}^\infty b_n \sin\frac{n\pi x}{L},
  \qquad
  b_n = \frac{2}{L}\int_0^L f(x)\sin\frac{n\pi x}{L}\,dx
\]

Both converge to $f(x)$ on $(0,L)$ at continuity points. At endpoints: cosine series converges to the actual endpoint values; sine series is forced to 0 (odd extension).

\subsection{Shift Property}

If $f$ has period $2\pi$ with coefficients $a_n,b_n$, and $g(x) = f(x+h)$, then:
\[
  A_n = a_n\cos(nh) + b_n\sin(nh),
  \qquad
  B_n = -a_n\sin(nh) + b_n\cos(nh)
\]
\textbf{Special case $h = \pi$:} $A_n = (-1)^n a_n$, $B_n = (-1)^n b_n$, $A_0 = a_0$.

\subsection{Parseval's Identity and Computing Numerical Series}

\[
  \frac{1}{L}\int_{-L}^{L}|f(x)|^2\,dx
  = \frac{a_0^2}{2} + \sum_{n=1}^\infty \left(a_n^2 + b_n^2\right)
\]

\textbf{Method 1 --- substitute $x = x_0$:} at a continuity point $S(x_0)=f(x_0)$; at a jump $S(x_0)=\frac{f(x_0^+)+f(x_0^-)}{2}$.
Useful values: $x=0$ (kills $\sin$ terms), $x=\pi$ ($\cos(n\pi)=(-1)^n$), $x=\pi/2$.

\textbf{Method 2 --- Parseval:} compute the left side as a definite integral; equate to the right side to extract the desired sum.

\textbf{Reference table:}
\begin{center}
\begin{tabular}{ll}
\toprule
Series & Value \\
\midrule
$\displaystyle\sum_{n=1}^\infty \dfrac{1}{n^2}$ & $\dfrac{\pi^2}{6}$ \\[10pt]
$\displaystyle\sum_{n=1}^\infty \dfrac{(-1)^{n+1}}{n^2}$ & $\dfrac{\pi^2}{12}$ \\[10pt]
$\displaystyle\sum_{n=0}^\infty \dfrac{1}{(2n+1)^2}$ & $\dfrac{\pi^2}{8}$ \\[10pt]
$\displaystyle\sum_{n=1}^\infty \dfrac{1}{n^4}$ & $\dfrac{\pi^4}{90}$ \\[10pt]
$\displaystyle\sum_{n=0}^\infty \dfrac{(-1)^n}{2n+1}$ & $\dfrac{\pi}{4}$ \\[10pt]
$\displaystyle\sum_{n=0}^\infty \dfrac{(-1)^n}{(2n+1)^3}$ & $\dfrac{\pi^3}{32}$ \\
\bottomrule
\end{tabular}
\end{center}

\begin{remark}
$\sum 1/(2n+1)^2 = \pi^2/8$ follows from $\sum 1/n^2 = \pi^2/6$ by separating even and odd terms:
$\sum_{n=1}^\infty \frac{1}{(2n)^2} = \frac{1}{4}\cdot\frac{\pi^2}{6}$, so
$\sum_{\text{odd}} = \frac{\pi^2}{6} - \frac{\pi^2}{24} = \frac{\pi^2}{8}$.
\end{remark}

\subsection{Complex Fourier Series}

\[
  f(x) = \sum_{n=-\infty}^\infty c_n\, e^{in\pi x/L},
  \qquad
  c_n = \frac{1}{2L}\int_{-L}^{L} f(x)\,e^{-in\pi x/L}\,dx
\]

\textbf{Relationship to real coefficients} ($n \ge 1$):
\[
  c_0 = \frac{a_0}{2},
  \qquad
  c_n = \frac{a_n - ib_n}{2},
  \qquad
  c_{-n} = \frac{a_n + ib_n}{2},
  \qquad
  a_n = c_n + c_{-n}, \quad b_n = i(c_n - c_{-n})
\]

\textbf{Quick route:} if you already have the real series, read off $c_n$ without recomputing integrals.

\textbf{Parseval for complex series:}
$\displaystyle\frac{1}{2L}\int_{-L}^L |f(x)|^2\,dx = \sum_{n=-\infty}^\infty |c_n|^2$

\subsection{Summing a Real Fourier Series in Closed Form}

\textbf{Geometric series technique:}
To find a closed form for $\displaystyle\sum_{n=1}^\infty r^n\sin(nx)$ or $\sum r^n\cos(nx)$ ($0<r<1$):
\begin{enumerate}
  \item Write as $\operatorname{Im}$ or $\operatorname{Re}$ of $\displaystyle\sum_{n=1}^\infty (re^{ix})^n$.
  \item Sum as a geometric series: $\displaystyle\sum_{n=1}^\infty z^n = \frac{z}{1-z}$,\quad $|z|<1$.
  \item Multiply numerator and denominator by $\overline{(1-z)}$; separate real and imaginary parts.
\end{enumerate}

\textbf{Example:}
\[
  \sum_{n=1}^\infty \frac{\sin(nx)}{3^n}
  = \operatorname{Im}\!\left(\frac{e^{ix}/3}{1-e^{ix}/3}\right)
  = \operatorname{Im}\!\left(\frac{3e^{ix}}{3-e^{ix}}\right)
  = \frac{3\sin x}{10-6\cos x}
\]
(denominator $|3-e^{ix}|^2 = 10-6\cos x$; imaginary part of numerator $= 3\sin x$).

\textbf{Infinite differentiability via the $M$-test on derivative series:}
$\sum u_n$ is $C^\infty$ on $D$ if for every $k\ge 0$ the series $\sum u_n^{(k)}$ passes the $M$-test
(i.e.\ $\sum_{n}\|u_n^{(k)}\|_\infty < \infty$).
\textbf{Example:} $|n^k\sin(nx)/3^n| \le n^k/3^n$, and $\sum n^k/3^n < \infty$ for all $k$ (geometric dominates polynomial),
so $\sum \sin(nx)/3^n \in C^\infty(\mathbb{R})$.

\subsection{Hyperbolic and Complex Functions (for Fourier)}

\textbf{Definitions and key identity:}
\[
  \sinh x = \frac{e^x - e^{-x}}{2}, \qquad \cosh x = \frac{e^x + e^{-x}}{2}, \qquad e^x = \cosh x + \sinh x
\]

\textbf{Properties:} $\cosh^2 x - \sinh^2 x = 1$;\quad $(\sinh x)' = \cosh x$;\quad $(\cosh x)' = \sinh x$;\\
$\sinh(-x) = -\sinh x$ (odd);\quad $\cosh(-x) = \cosh x$ (even).

\textbf{Euler's formula:} $e^{i\theta} = \cos\theta + i\sin\theta$

\textbf{All four functions via $e$:}
\begin{center}
\renewcommand{\arraystretch}{1.8}
\begin{tabular}{llll}
\toprule
$\cos\theta = \dfrac{e^{i\theta}+e^{-i\theta}}{2}$ &
$\sin\theta = \dfrac{e^{i\theta}-e^{-i\theta}}{2i}$ &
$\cosh x = \dfrac{e^{x}+e^{-x}}{2}$ &
$\sinh x = \dfrac{e^{x}-e^{-x}}{2}$ \\
\bottomrule
\end{tabular}
\end{center}

\textbf{Cross-relations} (multiplying argument by $i$ swaps trig $\leftrightarrow$ hyperbolic):
\[
  \cos(ix) = \cosh x, \quad \sin(ix) = i\sinh x, \quad \cosh(ix) = \cos x, \quad \sinh(ix) = i\sin x
\]

\textbf{Complex argument formulas:}
\[
  \sinh(a + ib) = \sinh(a)\cos(b) + i\cosh(a)\sin(b), \qquad
  \cosh(a + ib) = \cosh(a)\cos(b) + i\sinh(a)\sin(b)
\]

\textbf{Key simplification for integer $n$} (constant in complex Fourier):
\[
  \sinh(a - in\pi) = (-1)^n \sinh(a), \qquad \cosh(a - in\pi) = (-1)^n \cosh(a)
\]
because $\cos(n\pi) = (-1)^n$ and $\sin(n\pi) = 0$.

\subsection{Technique: Complex Fourier of $e^{ax}$ on $[-L,L]$}

\[
  c_n = \frac{1}{2L}\int_{-L}^L e^{ax}e^{-in\pi x/L}\,dx
      = \frac{(-1)^n \sinh(aL)}{L(a - in\pi/L)}
      = \frac{(-1)^n \sinh(aL)\,(a + in\pi/L)}{L(a^2 + n^2\pi^2/L^2)}
\]

Reading off the real coefficients:
\[
  a_n = \frac{2(-1)^n\sinh(aL)\cdot a}{L(a^2+n^2\pi^2/L^2)},
  \qquad
  b_n = \frac{2(-1)^n\sinh(aL)\cdot n\pi/L}{L(a^2+n^2\pi^2/L^2)}
\]

\textbf{Even/odd shortcut:} write $e^{ax} = \cosh(ax) + \sinh(ax)$.
Since $\cosh$ is even it contributes only to $a_n$; since $\sinh$ is odd it contributes only to $b_n$.
This halves the computation.

%----------------------------------------------------------------------
\section{Fourier Transform}

\subsection{Definition and Inverse ($\omega$ convention)}
\[
  \hat{f}(\omega) = \mathcal{F}\{f\}(\omega)
  = \int_{-\infty}^\infty f(x)\,e^{-i\omega x}\,dx,
  \qquad
  f(x) = \frac{1}{2\pi}\int_{-\infty}^\infty \hat{f}(\omega)\,e^{i\omega x}\,d\omega
\]

\textbf{Fourier Integral Representation} (real form):
\[
  f(x) = \int_0^\infty \bigl[A(\omega)\cos(\omega x) + B(\omega)\sin(\omega x)\bigr]\,d\omega,
  \quad
  A(\omega) = \frac{1}{\pi}\int_{-\infty}^\infty f(t)\cos(\omega t)\,dt,
  \quad
  B(\omega) = \frac{1}{\pi}\int_{-\infty}^\infty f(t)\sin(\omega t)\,dt
\]

\textbf{Pointwise convergence of the IFT:} if $f\in L^1(\mathbb{R})$ and $f$ is piecewise smooth,
the IFT formula holds at every point of continuity of $f$.

\textbf{Technique --- extracting definite integrals from the IFT:} substituting $x=x_0$ gives
\[
  \int_{-\infty}^\infty \hat{f}(\omega)\,e^{i\omega x_0}\,d\omega = 2\pi f(x_0).
\]
For even $f$ (so $\hat{f}$ is real): $\displaystyle\int_0^\infty \hat{f}(\omega)\cos(\omega x_0)\,d\omega = \pi f(x_0)$.

\subsection{Key Properties}

\begin{center}
\begin{tabular}{ll}
\toprule
\textbf{Property} & \textbf{Formula} \\
\midrule
Linearity          & $\mathcal{F}\{af+bg\} = a\hat{f}+b\hat{g}$ \\[4pt]
Time shift         & $\mathcal{F}\{f(x-a)\}(\omega) = e^{-i\omega a}\,\hat{f}(\omega)$ \\[4pt]
Freq.\ shift       & $\mathcal{F}\{e^{iax}f(x)\}(\omega) = \hat{f}(\omega - a)$ \\[4pt]
Scaling            & $\mathcal{F}\{f(ax)\}(\omega) = \dfrac{1}{|a|}\,\hat{f}\!\left(\dfrac{\omega}{a}\right)$ \\[4pt]
Derivative         & $\mathcal{F}\{f'(x)\}(\omega) = i\omega\,\hat{f}(\omega)$ \\[4pt]
Mult.\ by $x$      & $\mathcal{F}\{xf(x)\}(\omega) = i\,\dfrac{d\hat{f}}{d\omega}(\omega)$ \\[4pt]
Convolution        & $\mathcal{F}\{f*g\} = \hat{f}\cdot\hat{g}$ \\[4pt]
Duality            & $\mathcal{F}\{\hat{f}(x)\}(\omega) = 2\pi f(-\omega)$ \\[4pt]
Cos modulation     & $\mathcal{F}\{f(x)\cos(ax)\} = \dfrac{\hat{f}(\omega-a)+\hat{f}(\omega+a)}{2}$ \\[4pt]
Sin modulation     & $\mathcal{F}\{f(x)\sin(ax)\} = \dfrac{\hat{f}(\omega-a)-\hat{f}(\omega+a)}{2i}$ \\[4pt]
Even $f$           & $\hat{f}(\omega)$ is real \\[2pt]
Odd $f$            & $\hat{f}(\omega)$ is purely imaginary \\[4pt]
Riemann--Lebesgue  & If $f\in L^1(\mathbb{R})$: $\hat{f}$ is continuous on $\mathbb{R}$ and $\hat{f}(\omega)\to 0$ as $|\omega|\to\infty$ \\[4pt]
Half-line (even $f$) & $\displaystyle\int_0^\infty f(x)\cos(\omega x)\,dx = \tfrac{1}{2}\hat{f}(\omega)$
                       \quad (since $\hat{f}(\omega) = 2\int_0^\infty f(x)\cos(\omega x)\,dx$ when $f$ is even) \\
\bottomrule
\end{tabular}
\end{center}

\subsection{Common Transform Pairs}

\begin{center}
\begin{tabular}{lll}
\toprule
$f(x)$ & $\hat{f}(\omega)$ & Notes \\
\midrule
$e^{-ax}\,\mathbf{1}_{x\ge 0}$ & $\dfrac{1}{a+i\omega}$ & $a>0$ \\[10pt]
$e^{-a|x|}$ & $\dfrac{2a}{a^2+\omega^2}$ & $a>0$ \\[10pt]
$\dfrac{1}{x^2+a^2}$ & $\dfrac{\pi}{a}\,e^{-a|\omega|}$ & $a>0$ \\[10pt]
$e^{-ax^2}$ & $\sqrt{\dfrac{\pi}{a}}\,e^{-\omega^2/(4a)}$ & $a>0$, Gaussian $\to$ Gaussian \\[10pt]
$\mathbf{1}_{|x|\le L}$ & $\dfrac{2\sin(L\omega)}{\omega}$ & rect function \\
\bottomrule
\end{tabular}
\end{center}

\subsection{Parseval / Plancherel Theorem}
\[
  \int_{-\infty}^\infty |f(x)|^2\,dx
  = \frac{1}{2\pi}\int_{-\infty}^\infty |\hat{f}(\omega)|^2\,d\omega
\]

%----------------------------------------------------------------------
\section{Laplace Transform}

\subsection{Definition}
\[
  \mathcal{L}\{f\}(s) = F(s) = \int_0^\infty f(t)\,e^{-st}\,dt
\]

\subsection{Common Transforms and Properties}

\begin{center}
\begin{tabular}{ll}
\toprule
$f(t)$ & $F(s)$ \\
\midrule
$1$          & $\dfrac{1}{s}$ \\[8pt]
$t^n$        & $\dfrac{n!}{s^{n+1}}$ \\[8pt]
$e^{at}$     & $\dfrac{1}{s-a}$ \\[8pt]
$\sin(bt)$   & $\dfrac{b}{s^2+b^2}$ \\[8pt]
$\cos(bt)$   & $\dfrac{s}{s^2+b^2}$ \\[8pt]
$\sinh(bt)$  & $\dfrac{b}{s^2-b^2}$,\quad $s>|b|$ \\[8pt]
$\cosh(bt)$  & $\dfrac{s}{s^2-b^2}$,\quad $s>|b|$ \\[8pt]
$e^{at}f(t)$ & $F(s-a)$ \quad (frequency shift) \\[8pt]
$f'(t)$      & $sF(s)-f(0)$ \\[8pt]
$f''(t)$     & $s^2F(s)-sf(0)-f'(0)$ \\[8pt]
$(f*g)(t)$   & $F(s)\cdot G(s)$ \quad (convolution) \\
\bottomrule
\end{tabular}
\end{center}

\end{document}
